\chapter{Partial Completeness}
\label{ch:qtsl-partial-completeness}

Every rule in the catalogue of \bookref{ch:qtsl-manual}{the QTSL
manual} carries its own soundness story: a statement, a proof, and,
where mechanization has caught up, a machine-checked realization.
Soundness can therefore be audited locally. Each rule earns its keep
on its own, and a user who trusts the kernel need trust nothing else.

Completeness makes a different kind of claim, and no amount of local
auditing can establish it. A completeness theorem says that the rule
set as a whole is not missing anything, so that every true entailment
expressible in some class of specifications is actually derivable.
Claims of that shape must be won fragment by fragment, by exhibiting
either a derivation or a countermodel for every specification the
fragment can express. The prize is a change in the meaning of
refusal. Without completeness, a tool that fails to derive a
specification reports only that a rule might be missing; within a
complete fragment, the same failure delivers a definite verdict about
the specification itself.

This chapter records what has been proved. A reader who wants only
the shape of the result can stop after the completeness theorem
(Theorem~\ref{thm:gnand-completeness}); the axiom basis in
section~\ref{sec:gnand-proof-theory} shows what completeness costs,
and the closing sections account for the evidence and draw the
boundary of the claim. The most surprising idea in the chapter, for
a reader browsing, lives in Remark~\ref{rem:delayer-wings} and in
Figure~\ref{fig:delayer-wings}: the models that witness completeness
must branch, even though nothing in the specifications mentions
branching.

\section{A complete fragment: hard deadlines under staggering}
\label{sec:qtsl-gnand}

The fragment lives in the timed core of the calculus and uses two
modalities over propositional logic.

\begin{definition}[clocked model]
\label{def:clocked-model}
A \emph{clocked model} has finitely many locations, each carrying a
positive maximal dwell time and a fixed valuation of the atoms. A
\emph{trajectory} is a timed path through locations that respects the
dwell bounds.
\end{definition}

The two modalities read as follows.
\begin{itemize}
\item $\Box\varphi$, \emph{invariance}: $\varphi$ holds at every
  point reachable from here. The formula $D := \Box\bot$
  (``deadness'') holds exactly at the states with no future, so
  $\neg D$ asserts liveness.
\item $\mathsf{A}_T\varphi$, \emph{worst-case service}: \emph{every}
  trajectory from here reaches a $\varphi$-point within time at most
  $T$. The quantifier ranges over all branches, so the operator
  reports a worst case rather than a per-path promise. That one word,
  ``every'', is the source of everything difficult below.
\end{itemize}

The characteristic constraint of the fragment, and the origin of its
working name, is the \emph{gnand}\footnote{``Graded NAND'': a
prohibited conjunction, as in NAND, whose conjuncts are graded by
time bounds. The name arose as a working label during the
development and proved too convenient to retire.}, a globally held
prohibition of joint imminence:
\[
\mathrm{gnand}(a, d_a, b, d_b) \;:=\;
\Box\neg\big(\mathsf{A}_{d_a}a \wedge \mathsf{A}_{d_b}b\big),
\]
read: at no reachable moment is $a$ due within $d_a$ while $b$ is due
within $d_b$. It serves as a natural mutual-exclusion primitive for
services that must stay temporally staggered, two operations that may
never both be imminent, in either order.

\begin{example}[a single-boiler espresso machine]
\label{ex:espresso}
A single-boiler espresso machine cannot brew and steam at the same
time. Suppose an espresso order $e$ must be served within $90$
seconds on every branch of the schedule, a milk order $m$ within
$120$, and the boiler forbids any reachable moment at which both are
due within $30$ seconds. The specification comprises the pending
obligations $\neg e \wedge \mathsf{A}_{90}\,e$ and $\neg m \wedge
\mathsf{A}_{120}\,m$, the staggering constraint
$\mathrm{gnand}(e, 30, m, 30)$, and liveness $\neg D$. Suppose the
caf\'e also advertises that the $90$-second figure is tight, so that
some branch really does take the full $90$. That assertion,
$\neg\mathsf{A}_{89}\,e$, already forces the schedule to
\emph{branch}: a barista who always serves promptly would make the
worst case better than declared, so some branch of the schedule must
dawdle while another serves. This innocent observation grows into
the central construction of the completeness proof
(Remark~\ref{rem:delayer-wings}).
\end{example}

The example's shape, in general, defines the fragment.

\begin{definition}[gnand core]
\label{def:gnand-core}
A \emph{gnand core} is a specification asserting, over $n$ atoms:
\begin{itemize}
\item \textbf{pending obligations} $\neg a \wedge \mathsf{A}_{s_a}a$,
  meaning $a$ has not yet been served and must be served within $s_a$
  on every branch;
\item optionally, \textbf{exactness} $\neg\mathsf{A}_{s_a-1}a$ (the
  deadline is tight: some branch delays past $s_a - 1$) or an
  \textbf{avoidance fact} $\neg\mathsf{A}_h a$ (some branch withholds
  $a$ through time $h$);
\item any finite set of gnand constraints over pairs of atoms; and
  liveness $\neg D$.
\end{itemize}
Deadlines may be taken to be integers with no loss of
generality.\footnote{Rational rescaling and the decomposition of a
deadline into iterated unit deadlines are derivable rewrites in the
proof theory, so a fractional menu of deadlines reduces to an
integer one.}
\end{definition}

The fragment deliberately excludes triggered obligations
$\Box(x \to \mathsf{A}_d y)$, the atom-level variant of the
staggering constraint, and possibility assertions $\neg\Box\nu$ with
their associated time-locks. The boundary is principled rather than
incidental; it follows a \emph{sign}. Every constraint in the
fragment either forbids joint imminence or states a fact at the
initial state, and the model construction below is protected by a
monotonicity lemma that holds for such negative-side checks.
The excluded formers \emph{demand} positive worst-case facts at
interior states of a model. Constructions that meet such demands
exist in the development record, but their completeness theory is
open, and this section claims nothing about them.

\subsection{The proof theory}
\label{sec:gnand-proof-theory}

The axiom basis was not designed in advance; it was
\emph{discovered}, by iterating between attempted completeness proofs
and machine searches for countermodels. Three axiom schemata were
added to the seed basis this way, each certified genuinely new by an
explicit independence witness before being adopted, and several
others were pruned as derivable. The final basis comprises:
\begin{itemize}
\item four axioms governing the invariance modality;
\item reflexive service, $\varphi \to \mathsf{A}_T\varphi$;
\item budget composition,
  $\mathsf{A}_{T_1}\mathsf{A}_{T_2}\varphi \leftrightarrow
  \mathsf{A}_{T_1+T_2}\varphi$;
\item the graded absurdity law, $\mathsf{A}_T\bot \to D$;
\item relativization of obligations through invariants,
  $\Box\nu \wedge \mathsf{A}_T(\chi \wedge \nu) \leftrightarrow
  \Box\nu \wedge \mathsf{A}_T\chi$;
\item a first-service law: at the first service among $n$ co-pending
  obligations, the remaining obligations still bind; and
\item the last and decisive addition, a \emph{residual-split}
  first-service law. For any cutoff $c$ strictly below all deadlines,
  \[
  \bigwedge_{i \le n} \mathsf{A}_{T_i}\alpha_i \;\to\;
  \mathsf{A}_c\Big(\bigvee_{i}\big(\alpha_i \wedge
  \bigwedge_{i' \ne i} \mathsf{A}_{T_{i'}}\alpha_{i'}\big) \;\vee\;
  \bigwedge_i \mathsf{A}_{T_i - c}\alpha_i\Big),
  \]
  meaning that within $c$, either some first service has occurred
  (and the co-obligations persist), or nothing has been served and
  every obligation holds at its true residual $T_i - c$.
\end{itemize}
Derivation rules are propositional calculus, modus ponens, import of
a background theory, monotonicity for both modalities, and rational
rescaling.

\subsection{The completeness theorem}
\label{sec:gnand-theorem}

\begin{theorem}[completeness for the gnand fragment]
\label{thm:gnand-completeness}
For gnand cores, the following are equivalent:
\begin{enumerate}
\item the core is satisfiable in a clocked model;
\item an effective arithmetic on the core (a fixpoint computation on
  per-atom effective deadlines and service-precedence edges, with
  one additional branch-conditional check) fires none of its
  refutation conditions;
\item an explicit finite model of the core exists on the uniform
  grid with dwell $\tfrac{1}{n+1}$, where $n$ is the number of
  pending atoms, and the theorem's proof constructs it;
\item the core is consistent: $\not\vdash D$ from the core in the
  basis above.
\end{enumerate}
\end{theorem}

In particular, derivability and clocked validity coincide on the
fragment. Since soundness is proved over a much broader class of
trajectory spaces, of which clocked models are a special case,
validity over every intermediate semantic class coincides with both.
Satisfiability and validity of cores are therefore
\emph{decidable}, not merely semi-decidable; the fixpoint arithmetic
of item~(2) decides both.

Within its fragment the theorem is unconditional; its statement
carries no unproven hypotheses. Its proof, however, forms a chain
with several links, and the accounting must proceed link by link.
That every arithmetic refutation in item~(2) is matched by a formal
derivation of $D$ is a theorem, with the residual-split law doing
the work its adoption promised: the arithmetic's
deadline-tightening step is derivable by iterating the law at
successive cutoffs, descending step by step to the arithmetic's own
value. That every arithmetic-feasible core yields a model at grid
$\tfrac{1}{n+1}$ combines a proved monotone-protection lemma with one
combinatorial schema whose verification is per-instance and
population-scale (section~\ref{sec:gnand-evidence}) rather than
symbolic.

\begin{remark}[delayer wings]
\label{rem:delayer-wings}
The model construction deserves a name of its own. The emitted
models are \emph{schedule trees}: an earliest-deadline-first spine
together with \emph{delayer wings}, auxiliary branches whose only
role is to keep each declared worst case at its declared value, so
that the negative graded facts remain true even though every
individual path serves early (Figure~\ref{fig:delayer-wings}). That
branching structure can satisfy staggering constraints which every
branching-free model provably violates. This was the central
semantic discovery of the development, and it is why the fragment's
models are trees rather than runs.
\end{remark}

\begin{figure}[b]
\centering
\begin{tikzpicture}[
    >=Stealth,
    state/.style={draw, rounded corners=2pt, inner sep=4pt,
                  minimum height=2.6em,
                  font=\footnotesize, align=center},
    spine/.style={->, thick},
    wing/.style={->, thick, dashed, gray!55!black},
    lbl/.style={font=\footnotesize\itshape, midway, above=1pt},
    winglbl/.style={font=\footnotesize\itshape, midway, below=2pt,
                    sloped, gray!55!black}]
  \node[state] (root) {$\neg e,\ \neg m$\\ $e$ due $90$, $m$ due $120$};
  \node[state, right=30mm of root] (se) {$e$ served\\ $m$ due $75$};
  \node[state, right=16mm of se] (sm) {$m$ served};
  \node[state, below right=16mm and 6mm of root] (wait)
    {still $\neg e,\ \neg m$\\ $e$ due $30$, $m$ due $60$};
  \node[state, right=22mm of wait] (we) {$e$ served late\\ ($t = 90$)};
  \draw[spine] (root) -- node[lbl] {serve $e$ at $t=45$} (se);
  \draw[spine] (se) -- node[lbl] {serve $m$} (sm);
  \draw[wing] (root.south) -- node[winglbl] {dawdle to $t=60$} (wait.west);
  \draw[wing] (wait) -- node[winglbl] {serve $e$} (we);
\end{tikzpicture}
\caption{A schedule tree for the espresso core of
Example~\ref{ex:espresso}. The solid spine serves early, in
earliest-deadline-first order. The dashed delayer wing dawdles, so
that the worst-case service time of $e$ really attains the declared
$90$ seconds and the exactness assertion $\neg\mathsf{A}_{89}\,e$
comes out true. No single run can do both jobs; the tree does.}
\label{fig:delayer-wings}
\end{figure}

\subsection{Evidence and assurance status}
\label{sec:gnand-evidence}

The proofs are prose mathematics, developed against an unusually
dense instrument battery. The completeness chain itself is not yet
mechanized in a proof assistant. The three axiom schemata the search
added, however, now carry sorry-free machine-checked soundness
lemmas (\texttt{sound\_u10f}, \texttt{sound\_firstService}, and
\texttt{sound\_fsResidual}, in \texttt{Qtsl/TimedCore.lean} under
\texttt{lean/qtsl} in \path{maitria-coalition/qtsl}; their
statements are reproduced in Appendix~\ref{app:lean-cone}). A
machine-checked countermodel,
\texttt{firstService\_attainment\_necessary}, accompanies them. The
two first-service lemmas assume two properties of the trajectory space,
fusion closure and attainment of first service; clocked models
satisfy both. The countermodel shows that the attainment hypothesis
cannot be dropped, by exhibiting a fusion-closed space in which the
first-service schema's premises hold and its conclusion fails.
Rule-level mechanization otherwise tracks the
catalogue's own realization discipline (\bookref{ch:qtsl-manual}{the
QTSL manual}). The machine evidence behind the chain, all of it
replayable from the development record as it migrates to
\path{maitria-coalition/qtsl}:

\begin{center}
\begin{tabular}{@{}lll@{}}
\toprule
instrument & what it checks & result \\
\midrule
derivation sweep & arithmetic-refuted cores derive $D$ &
  $64{,}786$ cells, $0$ underivable \\
cross-pair sweep & the descent lemma's coverage &
  $4{,}116/4{,}116$ derived \\
branch engine & branch-conditional refutations derive $D$ &
  $1{,}000/1{,}000$; cycle mode $200/200$ \\
model emitter & feasible cores yield grid models &
  $122/122$ named; $541/541$ random \\
soundness controls & no false derivations of $D$ &
  silent on every satisfiable core \\
\bottomrule
\end{tabular}
\end{center}

The controls in the last row did real work during the development.
They caught two draft derivation engines asserting more than their
machinery guaranteed, both retired before any proof consumed them,
and the same battery falsified three earlier conjectured axiom bases
outright. A rule set that has survived that regime differs, as an
epistemic object, from one that has merely not yet failed.

\subsection{What is not claimed}
\label{sec:gnand-not-claimed}

Completeness is claimed for gnand cores only. The composition of
this fragment with triggered obligations and possibility constraints
is open theory; so is the calculus's next expressive tier, in which
worst-case service generalizes to a time-bounded \emph{until} that
carries safety clauses along the way. Minimality of the basis, a
separate question from sufficiency, is settled for integer deadline
menus (the residual-split law is genuinely underivable from the
rest: there is a six-state witness, and six states is minimal) and
open for arbitrarily fine fractional menus. No answer to the
minimality question can weaken the theorem, since a redundant axiom
would only mean that the same completeness holds from fewer
assumptions.
